\documentclass[twocolumn]{aastex631}
\usepackage{amsmath}
\usepackage{booktabs}
\shorttitle{Selection-aware SDSS BPT/sSFR study}
\shortauthors{NebulaMind}
\begin{document}

\title{Broad Optical BPT Galaxies and Catalog Specific Star Formation in SDSS DR17: A Selection-Aware Matched-Control Study}
\author{NebulaMind Research Autopilot}
\affiliation{Public SDSS DR17 data only}

\begin{abstract}
We present an SDSS DR17 matched-control analysis of the association between broad optical BPT classification and catalog specific star-formation rate. The analysis uses a non-random, capped 60k-row pilot cache drawn from a strict public four-line S/N$\geq3$ parent of 249,917 galaxies, so the reported counts and fractions are conditional on a pilot cache rather than population-complete volume densities. Broad BPT-selected galaxies are matched to star-forming controls in stellar mass and redshift only; the sample is not matched in morphology or aperture fraction, which can inflate the apparent offset if the BPT-selected population is more bulge-dominated, and broad low-ionization classes can also include LINER-like emission from retired stellar populations. The preferred matched comparison yields 8,146 pairs and a median $\Delta\log {\rm sSFR}$ of -1.309 dex, with a bootstrap interval of [-1.334,-1.283] dex. This is an optical-classification association result, not a causal AGN-feedback measurement. Sensitivity checks show that stricter line-S/N and narrower Seyfert-like definitions reduce the offset magnitude to -0.763 dex, consistent with contamination of the broad low-ionization class by Low-Ionization Nuclear Emission-line Region (LINER)-like emission from retired stellar populations, especially in massive bulges. Subclass and selection-function treatment must therefore precede any physical interpretation, and an accompanying supplementary denominator/proxy atlas collects the related baselines and missing-observable notes.
\end{abstract}

\keywords{galaxies: active --- galaxies: star formation --- galaxies: evolution --- surveys --- methods: statistical}

\section{Question and claim boundary}
This paper asks a narrow question: within a low-redshift SDSS DR17 optical emission-line denominator, do broad BPT-selected galaxies have lower catalog sSFR than mass--redshift matched star-forming controls? The answer is yes for the cached denominator analyzed here. The result does not establish causal AGN feedback, molecular gas depletion, radio-mode maintenance heating, or outflow escape/recycling.

The present scope also excludes morphology or aperture controls, Seyfert/LINER separation, AGN luminosity or Eddington-ratio proxies, gas-mass measurements, environment labels, and time-domain or duty-cycle modelling.

The claim boundary is part of the result. BPT line ratios classify optical excitation, not directly black-hole accretion power in every object; retired stellar populations and low-ionization nuclear emission-line region (LINER)-like ionization can contaminate broad low-ionization classes \citep{stasinska2008,stasinska2015}. For that reason the paper uses the phrase ``broad optical BPT AGN'' and treats stronger Seyfert-like cuts as a sensitivity check rather than as an interchangeable label.

\section{Data and shared selection}
The data backbone is public SDSS DR17 spectroscopy, photometry, emission-line measurements, and catalog physical-property estimates \citep{york2000,sdssdr17,brinchmann2004}. The pilot analysis sample is a capped 60k-row pilot cache selected sequentially by \texttt{specObjID} after an arbitrary pilot-query row limit; it is not a random sample. The strict public four-line S/N$\geq3$ eligible parent contains 249,917 rows, so the pilot cache covers 24.0\% of that strict parent.
Over the redshift interval $0.02<z<0.12$, the SDSS 3-arcsec fiber subtends roughly 1.2--6.5 kpc, so the catalog sSFR comparison is fiber-centered rather than global.
Because the 3-arcsec fiber samples only the central regions at low redshift, disk emission can be omitted and the catalog-derived total sSFR can be biased differently for bulge-dominated and disk-dominated systems.
The stellar-mass and sSFR values are taken from the public MPA-JHU-style value-added table \texttt{galSpecExtra}, using its catalog median estimators \texttt{lgm\_tot\_p50} and \texttt{specsfr\_tot\_p50} after joining \texttt{SpecObj}, \texttt{galSpecInfo}, and \texttt{PhotoObj}. Those are low-redshift SDSS catalog estimates, not rederived line-by-line physical measurements \citep{brinchmann2004,sdssdr17,york2000}.

\begin{deluxetable*}{lrrr}
\tabletypesize{\scriptsize}
\tablecaption{Selection cascade for the flagship denominator. The 60k-row pilot cache is an artificial pilot-query cap, not a physical selection effect.\label{tab:selection}}
\tablehead{\colhead{Selection stage} & \colhead{Public DR17 rows} & \colhead{Cached rows} & \colhead{Retention vs. spectro-z parent}}
\startdata
SpecObj GALAXY, 0.02<z<0.12 & 501,060 & -- & 100.0\% \\
plus galSpecInfo/PhotoObj/galSpecExtra and mass/sSFR bounds & 416,554 & -- & 83.1\% \\
plus galSpecLine join & 416,554 & -- & 83.1\% \\
four BPT lines positive with positive errors & 373,445 & 60,000 & 74.5\% \\
four BPT lines S/N>=3 & 249,917 & 60,000 & 49.9\% \\
four BPT lines S/N>=5 & 176,523 & 42,446 & 35.2\% \\
four BPT lines S/N>=10 & 91,768 & 22,311 & 18.3\% \\
\enddata
\tablecomments{Counts are read-only public SDSS DR17 count queries plus the cached analysis table. Retention is shown as a percentage of the spectro-z parent. Cached rows are shown only where the cache applies.}
\end{deluxetable*}

The selection is not neutral with respect to star formation. In public counts, S/N$\geq3$ in all four BPT lines keeps 33.6\% of the $-12<\log {\rm sSFR}<-11$ parent bin but 94.9\% of the $-10<\log {\rm sSFR}<-9.5$ bin. Marginal distribution checks between the pilot sample and the full public parent show no redshift, mass, or sSFR bin differing by more than 5 percentage points; the largest absolute differences are 2.03, -1.63, and -0.58 percentage points, respectively. That check is reassuring but does not remove the capped-cache limitation.

\section{Classification and matching}
BPT classes are computed from H$\alpha$, H$\beta$, [O~III]$\lambda5007$, and [N~II]$\lambda6584$ using standard demarcations \citep{baldwin1981,kewley2001,kauffmann2003bpt,kewley2006} (see Figure~\ref{fig:bpt}). The cached denominator contains 39,553 star-forming galaxies, 12,234 intermediate/composite galaxies, 8,146 broad optical AGN, and 67 unclassified objects. Each broad optical BPT galaxy is matched to the nearest star-forming control in standardized $(\log M_\star,z)$ space, with replacement, so the association still inherits any mismatch in structure or fiber coverage between the two populations. Matching is not performed in morphology, aperture fraction, halo mass, gas mass, AGN luminosity, or duty-cycle phase; these missing dimensions define follow-up requirements.

\begin{figure*}
\centering
\includegraphics[width=0.72\textwidth]{../figures/fig-bpt.pdf}
\caption{BPT line-ratio diagram for the cached SDSS DR17 denominator. The diagram verifies the optical-excitation classes used for matching; it does not by itself prove accretion-driven feedback.}
\label{fig:bpt}
\end{figure*}

\section{Matched-control result}
The preferred broad-BPT comparison gives a large negative catalog-sSFR offset for the broad BPT-selected galaxies relative to star-forming controls.
A median $\Delta\log {\rm sSFR}$ (target minus matched control) of -1.309 dex corresponds to roughly a 20-fold lower catalog sSFR within this fiber-centered matched comparison, but this manuscript does not convert that proxy offset into a global quenching threshold.
Because the comparison is fiber-centered and the matching ignores morphology, the measured offset is a relative difference between the broad-BPT and control samples within the fiber aperture; it could primarily reflect a comparison between bulge-dominated BPT-selected galaxies and disk-dominated star-forming controls, rather than a galaxy-wide suppression of star formation.

\begin{deluxetable*}{lrrrr}
\tabletypesize{\scriptsize}
\tablecaption{Robustness ladder for matched catalog-sSFR offsets.\label{tab:robust}}
\tablehead{\colhead{Variant} & \colhead{$N$ pairs} & \colhead{Median $\Delta\log {\rm sSFR}$} & \colhead{95\% interval} & \colhead{Interpretation}}
\startdata
Broad BPT AGN, S/N$\geq3$, nearest SF control with replacement & 8,146 & -1.309 & [-1.334,-1.283] & Preferred association estimate \\
Moderate mass--redshift caliper & 7,867 & -1.318 & -- & 96.6\% target coverage \\
Greedy no-replacement stress test & 7,419 & -1.446 & -- & Poorer balance; diagnostic only \\
Broad BPT AGN, S/N$\geq10$ & 1,530 & -0.744 & -- & Line-S/N sensitivity \\
N II Seyfert-like proxy, S/N$\geq3$ & 2,114 & -0.763 & -- & Subclass sensitivity; LINER-reduced proxy \\
\enddata
\tablecomments{$\Delta\log {\rm sSFR}$ is target minus matched star-forming control. The moderate mass--redshift caliper uses $|\Delta\log M_\star|\leq0.05$ and $|\Delta z|\leq0.002$. The Seyfert-like proxy excludes a portion of the LINER-like low-ionization tail by construction. All values are conditional on the optical emission-line denominator.}
\end{deluxetable*}

\begin{figure*}
\centering
\includegraphics[width=0.86\textwidth]{../figures/fig-matched-offsets.pdf}
\caption{Distribution of matched-pair catalog-sSFR offsets for broad optical BPT-selected galaxies minus nearest star-forming controls. The preferred estimate is strong within this denominator but changes under stricter line-S/N and narrower subclass definitions. The moderate matching caliper shown in Table~\ref{tab:robust} uses $|\Delta\log M_\star|\leq0.05$ and $|\Delta z|\leq0.002$.}
\label{fig:offsets}
\end{figure*}

\section{Interpretation}
The result is directly measured, reproducible, and falsifiable inside the stated denominator. The matched-offset distribution is shown in Figure~\ref{fig:offsets}. The median offset is large and survives a moderate mass--redshift caliper.

At the same time, the S/N$\geq10$ and Seyfert-like proxy variants reduce the magnitude from -1.309 dex to -0.763 dex (Table~\ref{tab:robust}), roughly half the preferred broad-BPT estimate. That behavior is consistent with the narrower proxy excluding a portion of the low-ionization tail, including low-ionization nuclear emission-line region (LINER)-like ionization from retired stellar populations and post-AGB stars in massive bulges, rather than identifying a different physical mechanism. The reduction in offset magnitude for stricter S/N and Seyfert-like subsets does not remove the morphology/aperture caveat: if the broad-BPT sample is more bulge-dominated than the star-forming controls, the -1.309 dex offset can be inflated relative to a global quenching signal. The broad contamination primarily affects the broad low-ionization selection, which is why the narrower Seyfert-like proxy yields the smaller offset. The most robust conclusion is therefore: broad optical BPT classification is associated with lower catalog sSFR in this capped 60k-row pilot cache. Claims about causal quenching require additional data: morphology and aperture controls, Seyfert/LINER separation, AGN luminosity or Eddington proxy, gas mass, environment, and time-domain/duty-cycle modelling.

\section{Conclusion}
RP-1 is best treated as a concise, selection-aware association paper. An accompanying supplementary denominator/proxy atlas holds the related topic notes until the missing observables are added. Those follow-up claims require radio, X-ray, CO/HI, resolved outflow, halo/group, or simulation-mock observables that are not present in the current SDSS-only analysis.
In practice, that means future work needs the kinds of measurements used in radio-mode, X-ray cavity, molecular-gas, outflow, environment, and simulation-mock studies \citep{best2005,dekel2006,fabian2012,heckmanbest2014,lamassa2013,mcnamara2007,veilleux2005,xcoldgass2017,xgass2018,cicone2014,carniani2017,fiore2017,simba2019,tng2019,eagle2015,peng2010,piotrowska2022,wetzel2013}; these references motivate the missing observables, but they are not part of the present SDSS-only denominator.

\section{Local reproducibility}
This PDF was generated from the local candidate package \texttt{RP1\_FLAGSHIP\_WITH\_SUPPLEMENT\_20260709T013510Z}. It does not replace any public-linked PDF and does not touch public pages, live roots, product databases, deployment state, git history, billing/OAuth state, cron jobs, or external submission systems.


\begin{thebibliography}{}
\bibitem[Abdurro'uf et al.(2022)]{sdssdr17} Abdurro'uf, Accetta, K., Aerts, C., et al. 2022, ApJS, 259, 35
\bibitem[Baldwin et al.(1981)]{baldwin1981} Baldwin, J.~A., Phillips, M.~M., \& Terlevich, R. 1981, PASP, 93, 5
\bibitem[Best et al.(2005)]{best2005} Best, P.~N., Kauffmann, G., Heckman, T.~M., et al. 2005, MNRAS, 362, 25
\bibitem[Brinchmann et al.(2004)]{brinchmann2004} Brinchmann, J., Charlot, S., White, S.~D.~M., et al. 2004, MNRAS, 351, 1151
\bibitem[Carniani et al.(2017)]{carniani2017} Carniani, S., Marconi, A., Maiolino, R., et al. 2017, A\&A, 605, A42
\bibitem[Catinella et al.(2018)]{xgass2018} Catinella, B., Saintonge, A., Janowiecki, S., et al. 2018, MNRAS, 476, 875
\bibitem[Cicone et al.(2014)]{cicone2014} Cicone, C., Maiolino, R., Sturm, E., et al. 2014, A\&A, 562, A21
\bibitem[Dave et al.(2019)]{simba2019} Dave, R., Angles-Alcazar, D., Narayanan, D., et al. 2019, MNRAS, 486, 2827
\bibitem[Dekel \& Birnboim(2006)]{dekel2006} Dekel, A., \& Birnboim, Y. 2006, MNRAS, 368, 2
\bibitem[Fabian(2012)]{fabian2012} Fabian, A.~C. 2012, ARA\&A, 50, 455
\bibitem[Fiore et al.(2017)]{fiore2017} Fiore, F., Feruglio, C., Shankar, F., et al. 2017, A\&A, 601, A143
\bibitem[Heckman \& Best(2014)]{heckmanbest2014} Heckman, T.~M., \& Best, P.~N. 2014, ARA\&A, 52, 589
\bibitem[Kauffmann et al.(2003a)]{kauffmann2003bpt} Kauffmann, G., Heckman, T.~M., Tremonti, C., et al. 2003a, MNRAS, 346, 1055
\bibitem[Kewley et al.(2001)]{kewley2001} Kewley, L.~J., Dopita, M.~A., Sutherland, R.~S., Heisler, C.~A., \& Trevena, J. 2001, ApJ, 556, 121
\bibitem[Kewley et al.(2006)]{kewley2006} Kewley, L.~J., Groves, B., Kauffmann, G., \& Heckman, T. 2006, MNRAS, 372, 961
\bibitem[LaMassa et al.(2013)]{lamassa2013} LaMassa, S.~M., Heckman, T.~M., Ptak, A., \& Urry, C.~M. 2013, ApJL, 765, L33
\bibitem[McNamara \& Nulsen(2007)]{mcnamara2007} McNamara, B.~R., \& Nulsen, P.~E.~J. 2007, ARA\&A, 45, 117
\bibitem[Nelson et al.(2019)]{tng2019} Nelson, D., Springel, V., Pillepich, A., et al. 2019, Computational Astrophysics and Cosmology, 6, 2
\bibitem[Peng et al.(2010)]{peng2010} Peng, Y.-j., Lilly, S.~J., Kovac, K., et al. 2010, ApJ, 721, 193
\bibitem[Piotrowska et al.(2022)]{piotrowska2022} Piotrowska, J.~M., Bluck, A.~F.~L., Maiolino, R., \& Peng, Y.-j. 2022, MNRAS, 512, 1052
\bibitem[Saintonge et al.(2017)]{xcoldgass2017} Saintonge, A., Catinella, B., Tacconi, L.~J., et al. 2017, ApJS, 233, 22
\bibitem[Schaye et al.(2015)]{eagle2015} Schaye, J., Crain, R.~A., Bower, R.~G., et al. 2015, MNRAS, 446, 521
\bibitem[Stasinska et al.(2008)]{stasinska2008} Stasinska, G., Asari, N.~V., Cid Fernandes, R., et al. 2008, MNRAS, 391, L29
\bibitem[Stasinska et al.(2015)]{stasinska2015} Stasinska, G., Costa Duarte, M.~V., Vale Asari, N., Cid Fernandes, R., \& Sodre, L. 2015, MNRAS, 449, 559
\bibitem[Veilleux et al.(2005)]{veilleux2005} Veilleux, S., Cecil, G., \& Bland-Hawthorn, J. 2005, ARA\&A, 43, 769
\bibitem[Wetzel et al.(2013)]{wetzel2013} Wetzel, A.~R., Tinker, J.~L., Conroy, C., \& van den Bosch, F.~C. 2013, MNRAS, 432, 336
\bibitem[York et al.(2000)]{york2000} York, D.~G., Adelman, J., Anderson, J.~E., Jr., et al. 2000, AJ, 120, 1579
\end{thebibliography}

\end{document}
